\documentclass[11pt]{article}
\usepackage{geometry}
\geometry{paperwidth=4.25in,paperheight=5in,margin=0.25in}
\usepackage{tikz}

%\usepackage[notextcomp]{kpfonts}
%\RequirePackage[T1]{fontenc}

\setlength{\parindent}{0in}

\begin{document}
\sffamily

\begin{center}
  \begin{tikzpicture}[line width=1pt,scale=0.6, every node/.style={transform shape}]
    \node[anchor=west] at (0.5in,0.2) {\scalebox{3.2}{\sffamily\Huge\textbf{XIMERA}}};
  \end{tikzpicture}
\end{center}

\vspace{-3ex}

\begin{tabular}{@{}c@{}c@{}}%
\begin{minipage}{0.5\textwidth}%
\begin{center}
\includegraphics[width=\textwidth]{ximera-qrcode.png}

\vspace{-4ex}
\scalebox{0.8}{\textsf{https://mooculus.osu.edu/}}%
\end{center}
\end{minipage} &%
\begin{minipage}{0.5\textwidth}%
\begin{center}
\huge
For Students! \\
For Teachers! \\
For Education Researchers! \\
\end{center}
\end{minipage} \\
\end{tabular}

\vfill

Building online courses isn't easy---but we want to change this with \textsl{Ximera}!

\vfill

\textsl{Ximera}, pronounced ``chimera,'' uses standard open-source
tools to build online content. All content is written in \LaTeX, a
standard markup-language and this content is stored on GitHub, an
open, online repository for text files. This means that
\textsl{Ximera} promotes both collaboration on content, and remixing
of content into new courses.



\end{document}

%%% Local Variables: 
%%% mode: latex
%%% TeX-master: t
%%% End: 
